\documentclass[10pt,a4paper]{amsart}
\usepackage[margin=19mm]{geometry}
\usepackage{amsmath,amssymb,mathtools}
\usepackage[scr=rsfs,cal=euler]{mathalfa}
\usepackage{xfrac}
\usepackage[unicode,hidelinks,bookmarks=false]{hyperref}
\newcommand{\ie}{i.e.\ }
\newcommand{\cf}{cf.\ }
\newcommand{\resp}{respectively\ }
\newcommand{\Subsection}{\subsection}
\providecommand{\nobreakdash}{\nobreak\hbox{-}}
\setlength{\emergencystretch}{2em}
\multlinegap0pt
\usepackage{tikz}
\usepackage{amssymb}
\usepackage{mathtools}
\usepackage{dsfont}
\newtheorem{proposition}{Proposition}
\newtheorem{lemma}{Lemma}
\newcommand{\Itwo}{\tilde{I}^{2}(n,d)}
\newcommand{\Pol}{\mathcal{P}}
\newcommand{\SdwreathA}{S_{d} \wreath A}
\newcommand{\bepsilon}{\boldsymbol{\varepsilon}}
\newcommand{\bi}{\mathbf{i}}
\newcommand{\bj}{\mathbf{j}}
\newcommand{\bone}{\mathbf{1}}
\newcommand{\bv}{\mathbf{v}}
\newcommand{\gl}{\mathfrak{gl}}
\newcommand{\parity}[1]{\overline{#1}}
\newcommand{\tI}{\tilde{I}}
\newcommand{\wreath}{\wr}
\title{Schur--Weyl Equivalences for Wreath Product Superalgebras}
\author{Lauren Grimley, Jonathan R. Kujawa}
\date{Source: arXiv:2509.18598v2 (CC BY 4.0)}
\begin{document}
\maketitle

\noindent\emph{Source and licence.} Schur--Weyl Equivalences for Wreath Product Superalgebras. Version arXiv:2509.18598v2 (CC BY 4.0), distributed by arXiv under the Creative Commons Attribution 4.0 International licence, http://creativecommons.org/licenses/by/4.0/, as recorded in the arXiv licence metadata for this exact version and shown on the abstract page https://arxiv.org/abs/2509.18598v2. Full licence notice: \texttt{licenses/arxiv-2509.18598-CC-BY-4.0.txt}.\par\smallskip

\noindent\textit{Editorial scope.} Counted unit: the article's lemma lemma (source0.tex lines 602--608). The article's complete original proof of it is delivered with it (source0.tex lines 609--664, one \texttt{proof} environment of 56 lines). No mathematics is written, repaired or re-derived: the statement and the proof are the article's own lines, in the article's own notation, and no retained line is substituted or re-flowed.  Retained uncounted as the source-local context the proof needs: lines 355--357; lines 517--519.  Nothing was omitted from any retained span, so the package carries no omission record.  No retained line contains a citation, so the wrapper carries no bibliography and no bibliography entry.  No retained line contains a figure environment, an image-inclusion command or drawing code, so no image is delivered and the package declares no asset dependency.  Editorial formatting only, in addition: an AMS \texttt{amsart} preamble that restates the article's own declarations and macros used by the retained lines, namely the environments \texttt{proposition}, \texttt{lemma} on separate counters, together with the standard packages those lines need.  The retained environments print in this document as Proposition 1, Lemma 1 in order of appearance, whereas the article prints the same items with the numbering of its own full document.  The complete native source of the article as distributed in the arXiv e-print for this exact version is bundled unchanged as \texttt{sources/185530/source0.tex}.
\begin{equation}\label{E:supercommutator}
[E_{r_{1},s_{1}}^{a_{1}}, E_{r_{2},s_{2}}^{a_{2}}] = \delta_{s_{1}, r_{2}} E_{r_{1},s_{2}}^{a_{1}a_{2}} - (-1)^{\parity{a_{1}}\cdot \parity{a_{2}}} \delta_{s_{2}, r_{1}} E_{r_{2},s_{1}}^{a_{2}a_{1}},
\end{equation}

\begin{proposition}\label{T:injective-theorem}
Let $n,d \geq 1$ with $n \geq d$.  Fix $\bi = (i_1,\dotsc , i_d) \in \tI(n,d)$.  Set $\bv^{\bone}_{\bi}=v^{1}_{i_1} \otimes \dotsb  \otimes v^{1}_{i_d}$.  Then the map $N \rightarrow V_{n}^{\otimes d}  \otimes_{\SdwreathA} N$ defined by $n \mapsto \bv^{\bone}_{\bi }  \otimes n$ is injective for any $\SdwreathA$-module $N$.
\end{proposition}

\begin{lemma}  The following statements are true:
\begin{enumerate}
\item If $(\bj, \bi)_{k}, (\bj', \bi)_{k} \in \Itwo$, then $\alpha^{a}_{(\bj , \bi)_{k}}=\alpha^{a}_{(\bj' , \bi)_{k}}$ for all homogeneous $a \in A$.
\item If $(\bj, \bi)_{k}, (\bj, \bi' )_{k} \in \Itwo$, then $\alpha^{a}_{(\bj , \bi)_{k}}=\alpha^{a}_{(\bj , \bi')_{k}}$ for all homogeneous $a \in A$.
\item If $(\bj, \bi)_{k} \in \Itwo$, then for any $p \in [1,n]\setminus \{i_{1}, \dotsc , i_{d}, j_{1}, \dotsc, j_{d} \}$ and $\ell \in [1,d]$ with $\ell \neq k$, if we define $\bi '=(i_{1}, \dotsc , i_{\ell-1}, p, i_{\ell +1}, \dotsc , i_{d})$ and $\bj'=(j_{1}, \dotsc , j_{\ell-1}, p, j_{\ell +1}, \dotsc , j_{d}) $, then $(\bj', \bi' )_{k} \in \Itwo$ and $\alpha^{a}_{(\bj , \bi)_{k}}=\alpha^{a}_{(\bj' , \bi')_{k}}$ for all homogeneous $a \in A$.
\end{enumerate}
\end{lemma}

\begin{proof}
To prove (1), we first consider the special case of the pair $(\bj, \bi)_{k}$ and $(\bi, \bi)_{k}$.  The claim in (1) is trivial in this case if $\bj = \bi$, so we assume $\bj  \neq \bi$ and, hence, $j_{k} \neq i_{k}$. Use the  supercommutator formula \eqref{E:supercommutator}  to see that
\[
E_{i_{k}, j_{k}}^{1}E_{j_{k}, i_{k}}^{a} = E_{j_{k},i_{k}}^{a}E_{i_{k},j_{k}}^{1} + E_{i_{k},i_{k}}^{a} - E^{a}_{j_{k}, j_{k}}.
\]
Letting both sides act on on $\bv^{}_{\bi} \otimes n$ yields 
\[
E_{i_{k}, j_{k}}^{1}E_{j_{k}, i_{k}}^{a}(\bv^{}_{\bi} \otimes n) = E_{i_{k}, j_{k}}^{1}(\bv^{}_{\bj} \otimes \alpha^{a}_{(\bj , \bi)_{k}}(n))= \bv^{}_{\bi} \otimes \alpha^{a}_{(\bj , \bi)_{k}}(n)
\] and 
\begin{align*}
\left(E_{j_{k},i_{k}}^{a}E_{i_{k},j_{k}}^{1} + E_{i_{k},i_{k}}^{a} - E^{a}_{j_{k}, j_{k}} \right)(\bv^{}_{\bi} \otimes n) &= E_{j_{k},i_{k}}^{a}E_{i_{k},j_{k}}^{1}(\bv^{}_{\bi} \otimes n) + E_{i_{k},i_{k}}^{a}(\bv^{}_{\bi} \otimes n) - E^{a}_{j_{k}, j_{k}}(\bv^{}_{\bi} \otimes n) \\
&= E_{j_{k},i_{k}}^{a}E_{i_{k},j_{k}}^{1}(\bv^{}_{\bi} \otimes n) + \bv^{}_{\bi} \otimes \alpha^{a}_{(\bi, \bi)_{k}}(n) - E^{a}_{j_{k}, j_{k}}(\bv^{}_{\bi} \otimes n).
\end{align*}
However, $\bepsilon_{\bi} + \varepsilon_{i_{k}}-\varepsilon_{j_{k}} \not \in \Lambda(n,d)$ and so the first term of the final sum equals zero; likewise, the observation that $h_{j_{k}}\bv_{\bi}=0$ along with the second property of supermodules in $\Pol_{d}(\gl_{n}(A))$ implies the last term is also equal to zero.  Therefore, $\bv^{}_{\bi} \otimes \alpha^{a}_{(\bj , \bi)_{k}}(n) = \bv^{}_{\bi} \otimes \alpha^{a}_{(\bi, \bi)_{k}}(n)$.  It follows from Proposition~\ref{T:injective-theorem} that $\alpha^{a}_{(\bj , \bi)_{k}}(n) = \alpha^{a}_{(\bi, \bi)_{k}}(n)$, as desired.

Applying the special case to the pair $(\bj , \bi )_{k}$ and $(\bi,\bi)_{k}$, and the pair $(\bj' , \bi )_{k}$ and $(\bi,\bi)_{k}$, yields $\alpha^{a}_{(\bj , \bi)_{k}} = \alpha^{a}_{(\bi,\bi)_{k}}$ and $ \alpha^{a}_{(\bj', \bi)}= \alpha^{a}_{(\bi,\bi)_{k}}$, respectively.  Combining the two gives  $\alpha^{a}_{(\bj , \bi)_{k}}  = \alpha^{a}_{(\bj', \bi)}$, as desired.


Likewise, to prove (2), it suffices to consider the case of $(\bj, \bi)_{k}$ and $(\bj, \bj )_{k}$ where $\bj \neq \bi$.  In particular, $j_{k} \neq i_{k}$.  Use the supercommutator formula \eqref{E:supercommutator} to see that 
\[
E^{a}_{j_{k}, j_{k}}E^{1}_{j_{k}, i_{k}} = E^{1}_{j_{k}, i_{k}}E^{a}_{j_{k}, j_{k}} + E^{a}_{j_{k}, i_{k}}
\]  Applying each side to $ \bv_{\bi}^{} \otimes n$ yields
\[
E^{a}_{j_{k}, j_{k}}E^{1}_{j_{k}, i_{k}}( \bv_{\bi}^{} \otimes n) = E^{a}_{j_{k}, j_{k}}\left( \bv_{\bj}^{} \otimes n \right) =  \bv_{\bj}^{} \otimes \alpha_{(\bj , \bj)_{k}}^{a} (n)
\]
and 
\begin{align*}
\left( E^{1}_{j_{k}, i_{k}}E^{a}_{j_{k}, j_{k}} + E^{a}_{j_{k}, i_{k}}\right)( \bv_{\bi} \otimes n) &= \left( E^{1}_{j_{k}, i_{k}}E^{a}_{j_{k}, j_{k}}\right)( \bv_{\bi} \otimes n) +  \bv_{\bj} \otimes  \alpha_{(\bj, \bi)_{k}}^{a}(n).
\end{align*}  However, the first term in the last sum is again zero by the second property of supermodules in $\Pol_{d}(\gl_{n}(A))$.  Proposition~\ref{T:injective-theorem} then implies $\alpha_{(\bj , \bj)_{k}}^{a} (n) = \alpha_{(\bj, \bi)_{k}}^{a}(n)$, as claimed.


\iffalse
To prove (2), by symmetry it we may assume $(\bj',\bi)_{k} \in \Itwo$.   We also note that if $\bi = \bi'$, then the claim is covered by (1). So we may also assume $\bi \neq \bi'$ and, in particular, $i_{k} \neq i'_{k}$. We use this along with the fact that $j_{k}\neq i_{k}$, $j'_{k} \neq i'_{k}$, and $j'_{k} \neq i_{k}$.  Use the supercommutator formula \eqref{E:supercommutator} to see that 
\[
E^{a}_{j'_{k},i'_{k}}E_{i'_{k},i_{k}}^{1} =E_{i'_{k},i_{k}}^{1}E^{a}_{j'_{k},i'_{k}} + E_{j'_{k},i_{k}}^{a}.
\]  Applying each side to $ \bv_{\bi}^{} \otimes n$ yields
\[
E^{a}_{j'_{k},i'_{k}}E_{i'_{k},i_{k}}^{1}( \bv_{\bv_{\bi}}^{}\otimes n) = E^{a}_{j'_{k},i'_{k}}( \bv_{\bv_{\bi'}}^{} \otimes n) =  \bv_{\bj'}^{} \otimes \alpha_{(\bj' , \bi')_{k}}^{a}(n) 
\]
and 
\begin{align*}
\left( E_{i'_{k},i_{k}}^{1}E^{a}_{j'_{k},i'_{k}} + E_{j'_{k},i_{k}}^{a}\right)( \bv_{\bi}\otimes n) &= \left( E_{i'_{k},i_{k}}^{1}E^{a}_{j'_{k},i'_{k}}\right)( \bv_{\bi} \otimes n) +  \bv_{\bj'}\otimes \alpha_{(\bj', \bi)_{k}}^{a}(n).
\end{align*}  However, $\varepsilon_{\bi} + \varepsilon_{j'_{k}} - \varepsilon_{i'_{k}}$ is not a polynomial weight and so the first term equals zero.  Moreover, by weight considerations we see that the third term also equals zero.  Injectivity of the map $n' \mapsto  \bv_{\bj'} \otimes n'$  implies $\alpha_{(\bj' , \bi')_{k}}^{a} (n) = \alpha_{(\bj', \bi)_{k}}^{a}(n)$.  But (1) implies $\alpha_{(\bj', \bi)}^{a}(n)=\alpha_{(\bj, \bi)}^{a}(n)$.  Combining these two equalities yields the claim.
\fi

To prove (3), we first note that, by the choice of $p$, the entries of $\bi'$ and $\bj'$ are pairwise distinct and the two have the same values in every position except possibly at $i'_{k}$ and $j'_{k}$.  Therefore $(\bi', \bj')_{k} \in \Itwo$, as claimed.  Noting that $j'_{k}=j_{k}$ and $i'_{k}=i_{k}$, and, hence, that $i'_{k}\neq p$ and $j'_{k}\neq i_{\ell}$, the supercommutator formula \eqref{E:supercommutator} implies that 
\[
E^{a}_{j'_{k},i'_{k}}E_{p,i_{\ell}}^{1} =E_{p,i_{\ell}}^{1}E^{a}_{j'_{k},i'_{k}}.
\] Applying each side to $ \bv_{\bi}^{} \otimes n$ yields
\[
\left(E^{a}_{j'_{k},i'_{k}}E_{p,i_{\ell}}^{1} \right)( \bv_{\bi}^{} \otimes n) = E^{a}_{j'_{k},i'_{k}}( \bv_{\bi'}^{} \otimes n)= \bv_{\bj'}^{} \otimes \alpha_{(\bj', \bi')_{k}}^{a}(n) 
\] and 
\[
E_{p,i_{\ell}}^{1}E^{a}_{j'_{k},i'_{k}}( \bv_{\bi}^{} \otimes n) = E_{p,i_{\ell}}^{1}E^{a}_{j_{k},i_{k}}( \bv_{\bi}^{} \otimes n)=E_{p,j_{\ell}}^{1}\left( \bv_{\bj}^{} \otimes \alpha_{\bj, \bi}^{a} (n) \right) =  \bv_{\bj'}^{} \otimes \alpha_{\bj, \bi}^{a} (n).
\]  Proposition~\ref{T:injective-theorem} implies $\alpha_{(\bj' , \bi')_{k}}^{a} (n) = \alpha_{(\bj, \bi)_{k}}^{a}(n)$, as claimed.
\end{proof}

\end{document}
